\documentclass[a4paper,10pt]{article}
\usepackage[utf8]{inputenc}
\usepackage[english]{babel}
\usepackage{amsmath, amssymb, mathtools}
\usepackage{tikz}
\usetikzlibrary{patterns}
\usepackage{geometry}
\geometry{margin=1.7cm}
\usepackage{hyperref}
\usepackage{enumitem}
\usepackage{nicematrix}
\usepackage{eurosym} % para el euro
\setlength{\parindent}{0pt}
\title{Linear Programming\\Algorithmic Methods for Mathematical Models}
\author{Elena de Paz Obiols}
\date{\today}
\usepackage{float}

\begin{document}



\maketitle

\section*{Problem 1: Workstations}

A company makes three products and has available four workstations. The production time (in
minutes) per unit produced varies from workstation to workstation as shown below:

\begin{table}[H]
\centering
\begin{tabular}{|c|c|c|c|c|}
\hline
& Workstation 1 & Workstation 2 & Workstation 3 & Workstation 4 \\
\hline
Product 1 & 5 & 7 & 4 & 10 \\
\hline
Product 2 & 6 & 12 & 8 & 15 \\
\hline
Product 3 & 13 & 14 & 9 & 17 \\
\hline
\end{tabular}
\end{table}

Similarly the profit (e) contribution per unit varies from workstation to workstation as below:

\begin{table}[H]
\centering
\begin{tabular}{|c|c|c|c|c|}
\hline
& Workstation 1 & Workstation 2 & Workstation 3 & Workstation 4 \\
\hline
Product 1 & 10 & 8 & 6 & 9 \\
\hline
Product 2 & 18 & 20 & 15 & 17 \\
\hline
Product 3 & 15 & 16 & 13 & 17 \\
\hline
\end{tabular}
\end{table}

If next week there are 35 working hours available at each workstation, how much of each product
should be produced to maximize profit given that we need at least 100 units of product 1, 150
units of product 2 and 100 units of product 3?

Formulate this problem as a linear program.


\medskip
\textbf{Solution:} 

\textbf{Decision variables:}

\begin{itemize}
    \item $x_{ij}$: number of units of product $i$ produced at workstation $j$, for $i=1,2,3$ and $j=1,2,3,4$.
\end{itemize}

\textbf{Objective function:}
\begin{equation*}
    \text{Maximize } Z = 10x_{11} + 8x_{12} + 6x_{13} + 9x_{14} + 18x_{21} + 20x_{22} + 15x_{23} + 17x_{24} + 15x_{31} + 16x_{32} + 13x_{33} + 17x_{34} 
\end{equation*}
\textbf{Constraints:}
\begin{enumerate}
    \item Time constraints at each workstation:
    \begin{align*}
        5x_{11} + 6x_{21} + 13x_{31} &\leq 2100 \quad \text{(Workstation 1)} \\
        7x_{12} + 12x_{22} + 14x_{32} &\leq 2100 \quad \text{(Workstation 2)} \\
        4x_{13} + 8x_{23} + 9x_{33} &\leq 2100 \quad \text{(Workstation 3)} \\
        10x_{14} + 15x_{24} + 17x_{34} &\leq 2100 \quad \text{(Workstation 4)}
    \end{align*}
    
    \item Minimum production requirements:
    \begin{align*}
        x_{11} + x_{12} + x_{13} + x_{14} &\geq 100 \quad \text{(Product 1)} \\
        x_{21} + x_{22} + x_{23} + x_{24} &\geq 150 \quad \text{(Product 2)} \\
        x_{31} + x_{32} + x_{33} + x_{34} &\geq 100 \quad \text{(Product 3)}
    \end{align*}
    
    \item Non-negativity constraints:
    \begin{align*}
        x_{ij} &\geq 0 \quad \text{for all } i=1,2,3; j=1,2,3,4
    \end{align*}
\end{enumerate}

\newpage
\section*{Problem 2: Cargo plane}

A cargo plane has three compartments for storing cargo: front, centre and rear. These compart
ments have the following limits on both weight and volume:

\begin{table}[H]
\centering
\begin{tabular}{|c|c|c|}
\hline
Compartment & Weight capacity (tonnes) & Volume capacity (m$^3$) \\
\hline
Front & 10 & 6800 \\
\hline
Centre & 16 & 8700 \\
\hline
Rear & 8 & 5300 \\
\hline
\end{tabular}
\end{table}

Furthermore, for each compartment the ratio (weight in that compartment) / (weight capacity
in that compartment) must be the same, to maintain the balance of the plane.
The following four cargoes are available for shipment on the next flight:

\begin{table}[H]
\centering
\begin{tabular}{|c|c|c|c|}
\hline
Cargo & Weight (tonnes) & Volume (m$^3$/tonne) & Profit ($\euro$/tonne) \\
\hline
C1 & 18 & 480 & 310 \\
\hline
C2 & 15 & 650 & 380 \\
\hline
C3 & 23 & 580 & 350 \\
\hline
C4 & 12 & 390 & 285 \\
\hline
\end{tabular}
\end{table}

ny proportion of these cargoes can be accepted. The objective is to determine how much (if
any) of each cargo C1, C2, C3 and C4 should be accepted and how to distribute each among
the compartments so that the total profit for the flight is maximised.


Formulate the above problem as a linear program.

\medskip
\textbf{Solution:}

\textbf{Decision variables:}
\begin{itemize}
    \item $x_{ij}$: amount of cargo $i$ stored in compartment $j$, for $i=1,2,3,4$ and $j=F,C,R$ (Front, Centre, Rear).
\end{itemize}
\textbf{Objective function:}
\begin{equation*}
    \text{Maximize } Z = 310(x_{1F} + x_{1C} + x_{1R}) + 380(x_{2F} + x_{2C} + x_{2R}) + 350(x_{3F} + x_{3C} + x_{3R}) + 285(x_{4F} + x_{4C} + x_{4R})
\end{equation*}
\textbf{Constraints:}
\begin{enumerate}
    \item Weight capacity constraints:
    \begin{align*}
        x_{1F} + x_{2F} + x_{3F} + x_{4F} &\leq 10 \quad \text{(Front)} \\
        x_{1C} + x_{2C} + x_{3C} + x_{4C} &\leq 16 \quad \text{(Centre)} \\
        x_{1R} + x_{2R} + x_{3R} + x_{4R} &\leq 8 \quad \text{(Rear)}
    \end{align*}
    
    \item Volume capacity constraints:
    \begin{align*}
        480x_{1F} + 650x_{2F} + 580x_{3F} + 390x_{4F} &\leq 6800 \quad \text{(Front)} \\
        480x_{1C} + 650x_{2C} + 580x_{3C} + 390x_{4C} &\leq 8700 \quad \text{(Centre)} \\
        480x_{1R} + 650x_{2R} + 580x_{3R} + 390x_{4R} &\leq 5300 \quad \text{(Rear)}
    \end{align*}
    
    \item Balance constraint:
    \begin{align*}
        \frac{x_{1F} + x_{2F} + x_{3F} + x_{4F}}{10} = \frac{x_{1C} + x_{2C} + x_{3C} + x_{4C}}{16} = \frac{x_{1R} + x_{2R} + x_{3R} + x_{4R}}{8}
    \end{align*}
    
    \item Non-negativity constraints:
    \begin{align*}
        x_{ij} &\geq 0 \quad \text{for all } i=1,2,3,4; j=F,C,R
    \end{align*}
\end{enumerate}

\newpage
\section*{Problem 3: Canning}

 A canning company operates two canning plants. The growers are willing to supply fresh fruits
in the following amounts:
\begin{itemize}
    \item S1: 200 tonnes at 11 $\euro$/tonne
    \item S2: 310 tonnes at 10 $\euro$/tonne
    \item S3: 420 tonnes at 9 $\euro$/tonne
\end{itemize}

Shipping costs in $\euro$ per tonne are:

\begin{table}[H]
\centering
\begin{tabular}{|c|c|c|}
\hline
& Plant A & Plant B \\
\hline
S1 & 3 & 3.5 \\
\hline
S2 & 2 & 2.5 \\
\hline
S3 & 6 & 4 \\
\hline
\end{tabular}
\end{table}

Plant capacities and labour costs are:

\begin{table}[H]
\centering
\begin{tabular}{|c|c|c|}
\hline
& Plant A & Plant B \\
\hline
Capacity & 460 tonnes & 560 tonnes \\
\hline
Labour cost & 26 \euro/tonne & 21 \euro/tonne \\
\hline
\end{tabular}
\end{table}

The canned fruits are sold at 50 e/tonne to the distributors. The company can sell at this price
all they can produce. The objective is to find the best mixture of the quantities supplied by the
three growers to the two plants so that the company maximises its profits.


Formulate the above problem as a linear program.

\medskip
\textbf{Solution:}

\textbf{Decision variables:}
\begin{itemize}
    \item $x_{ij}$: amount of fruit supplied by grower $i$ to plant $j$, for $i=1,2,3$ and $j=A,B$.
\end{itemize}
\textbf{Objective function:}
\begin{equation*}
    \text{Maximize } Z = 50(x_{1A} + x_{1B} + x_{2A} + x_{2B} + x_{3A} + x_{3B}) - (11x_{1A} + 10x_{2A} + 9x_{3A} + 11x_{1B} + 10x_{2B} + 9x_{3B}) - (3x_{1A} + 2x_{2A} + 6x_{3A} + 3.5x_{1B} + 2.5x_{2B} + 4x_{3B}) - (26(x_{1A} + x_{2A} + x_{3A}) + 21(x_{1B} + x_{2B} + x_{3B}))
\end{equation*}
\textbf{Constraints:}
\begin{enumerate}
    \item Supply constraints:
    \begin{align*}
        x_{1A} + x_{1B} &\leq 200 \quad \text{(Grower S1)} \\
        x_{2A} + x_{2B} &\leq 310 \quad \text{(Grower S2)} \\
        x_{3A} + x_{3B} &\leq 420 \quad \text{(Grower S3)}
    \end{align*}
    
    \item Plant capacity constraints:
    \begin{align*}
        x_{1A} + x_{2A} + x_{3A} &\leq 460 \quad \text{(Plant A)} \\
        x_{1B} + x_{2B} + x_{3B} &\leq 560 \quad \text{(Plant B)}
    \end{align*}
    
    \item Non-negativity constraints:
    \begin{align*}
        x_{ij} &\geq 0 \quad \text{for all } i=1,2,3; j=A,B
    \end{align*}
\end{enumerate}

\newpage
\section*{Problem 4: Manufacturer}
A company manufactures four products (1, 2, 3, 4) on two machines (A and B). The time (in
minutes) to process one unit of each product on each machine is shown below:

\begin{table}[H]
\centering
\begin{tabular}{|c|c|c|}
\hline
& Machine A & Machine B \\
\hline
Product 1 & 10 & 27 \\
\hline
Product 2 & 12 & 19 \\
\hline
Product 3 & 13 & 33 \\
\hline
Product 4 & 8 & 23 \\
\hline
\end{tabular}
\end{table}

The profit per unit for each product (1, 2, 3, 4) is 10 \euro, 12 \euro, 17 \euro and 8 \euro respectively.
Each unit of product 1 needs to be processed by both machines A and B (let’s say first A and
then B, although this detail is not important). On the other hand, products 2, 3 and 4 can be
produced with a single machine, A or B.

The factory is very small and this means that floor space is very limited. Only one week’s
production is stored in 50 square metres of floor space where the floor space taken up by one unit
of each product is 0.1, 0.15, 0.5 and 0.05 (square metres) for products 1, 2, 3 and 4 respectively.

Due to customer requirements, the amount of product 3 that is produced should be related to
the amount of product 2 that is produced. Namely, over a week approximately the number of
units of product 2 to be produced should be twice as many units of product 3 (±5%).

Machine A is out of action (for maintenance/because of breakdown) 5$\%$ of the time and machine
B 7$\%$ of the time.

Assuming a working week 35 hours long formulate the problem of how to manufacture these
products.

\medskip
\textbf{Solution:}

\textbf{Decision variables:}
\begin{itemize}
    \item $x_{1}$: number of units of product 1 produced.
    \item $x_{p,m}$: number of units of product $p$ produced on machine $m$, for $p=2,3,4$ and $m=A,B$.
\end{itemize}
\textbf{Objective function:}
\begin{equation*}
    \text{Maximize } Z = 10x_{1} + 12(x_{2A} + x_{2B}) + 17(x_{3A} + x_{3B}) + 8(x_{4A} + x_{4B})
\end{equation*}
\textbf{Constraints:}
\begin{enumerate}
    \item Machine time constraints:
    \begin{align*}
        10x_{1} + 12x_{2A} + 13x_{3A} + 8x_{4A} &\leq 1995 \quad \text{(Machine A)} \\
        27x_{1} + 19x_{2B} + 33x_{3B} + 23x_{4B} &\leq 2265 \quad \text{(Machine B)}
    \end{align*}
    
    \item Floor space constraint:
    \begin{align*}
        0.1x_{1} + 0.15(x_{2A} + x_{2B}) + 0.5(x_{3A} + x_{3B}) + 0.05(x_{4A} + x_{4B}) &\leq 50
    \end{align*}
    
    \item Product relationship constraint:
    \begin{align*}
        2(x_{3A} + x_{3B}) - (x_{2A} + x_{2B}) &\leq 0.05(x_{2A} + x_{2B}) \\
        (x_{2A} + x_{2B}) - 2(x_{3A} + x_{3B}) &\leq 0.05(x_{2A} + x_{2B})
    \end{align*}
    
    \item Non-negativity constraints:
    \begin{align*}
        x_{1}, x_{p,m} &\geq 0 \quad \text{for all } p=2,3,4; m=A,B
    \end{align*}
\end{enumerate}

\end{document}